\documentclass[aspectratio=169]{beamer}

% ===== 基本パッケージ =====
\usepackage{adjustbox}
\usepackage{appendixnumberbeamer}
\usepackage[english]{babel}
\usepackage[utf8x]{inputenc}
\usepackage{amssymb,booktabs,amsfonts,mathtools}
\usepackage[babel]{csquotes}
\usepackage{subfiles}
\usepackage{threeparttable}
\usepackage{graphicx}
\usepackage[round]{natbib}
\usepackage{tabularray}
\usepackage{float}
\usepackage[normalem]{ulem}
\usepackage{multirow}
\usepackage{multicol}
\usepackage{comment}
\usepackage{tikz}
\usetikzlibrary{arrows.meta}

% ===== フォント =====
\usefonttheme{professionalfonts}
\renewcommand{\familydefault}{\sfdefault}

% ===== カラー定義 =====
\definecolor{myOrange}{RGB}{230,120,0}
\definecolor{myDarkGray}{RGB}{50,50,50}

% ===== カラーテーマ設定 =====
\setbeamercolor{structure}{fg=myOrange}
\setbeamercolor{title}{fg=myOrange}
\setbeamercolor{frametitle}{fg=white, bg=myOrange}
\setbeamercolor{normal text}{fg=myDarkGray, bg=white}

% ===== フレームタイトル =====
\setbeamertemplate{frametitle}{
	\vspace{0.2em}
	\begin{beamercolorbox}[wd=\paperwidth,ht=2.8ex,dp=1.2ex,leftskip=1em]{frametitle}
		\insertframetitle
	\end{beamercolorbox}
}

% ===== フッター（左：著者 / 中：スライドタイトル / 右：ページ）=====
\setbeamertemplate{footline}{
	\leavevmode
	\hbox{
		% 左
		\begin{beamercolorbox}[wd=.33\paperwidth,ht=2.5ex,dp=1.2ex,left]{author in head/foot}
			\hspace{1em}\insertshortauthor
		\end{beamercolorbox}
		% 中
 		 \begin{beamercolorbox}[wd=.33\paperwidth,ht=2.5ex,dp=1.2ex,center]{title in head/foot}% \hspace{1em}
			\usebeamerfont{title in head/foot}
		\end{beamercolorbox}%
		% 右
		\begin{beamercolorbox}[wd=.33\paperwidth,ht=2.5ex,dp=1.2ex,right]{date in head/foot}
			\insertframenumber{} / \inserttotalframenumber\hspace{1em}
		\end{beamercolorbox}
	}
}

\setbeamercolor{author in head/foot}{fg=white,bg=myOrange}
\setbeamercolor{title in head/foot}{fg=white,bg=myOrange}
\setbeamercolor{date in head/foot}{fg=white,bg=myOrange}

% ===== ナビゲーション消去 =====
\setbeamertemplate{navigation symbols}{}

% ===== 情報 =====
\makeatletter
\newcommand*\exInput[1]{\@@input#1 }
\makeatother

\title[Short Title]{Results of Baseline and Bi-weekly Surveys}
\author[Takahashi et al.]{Kazushi Takahashi, Wataru Kodama, Yukichi Mano}
\date{\today}
\begin{document}
	
	% --- 1枚目のスライド（タイトルページ） ---
	\begin{frame}{Title}
		\titlepage
	\end{frame}
	
	% --- 2枚目のスライド ---
	\begin{frame}{Contents} % 
		\begin{itemize}
			\item Motivation % 
			\item Research Questions and What We Do
			\item Study Areas and Research Design
			\item Descriptive statistics
			\item Preliminary results 
			\item Ways forward  
		\end{itemize}
	\end{frame}
% --- Motivation ---
\begin{frame}{Motivation}
	
	\begin{itemize}
		\item Improving agricultural productivity requires the adoption of new technologies,
		such as improved seeds, chemical fertilizers, and appropriate crop management practices.
		
		\vspace{0.3cm}
		
		\item Since technology development, adoption, and diffusion involve externality,
		public institutions play a central role.
		Public agricultural extension workers are key actors in disseminating technologies
		developed at research stations.
		
		\vspace{0.3cm}
		
		\item However, in many developing countries, particularly in Sub-Saharan Africa,
		public extension systems often function poorly.
		
		\vspace{0.3cm}
		
		\item One explanation is weak incentives:
		\begin{itemize}
			\item Limited monitoring and weak linkage between performance and rewards
			(lack of extrinsic motivation)
			\item Extension workers may not perceive extension activities as meaningful
			(lack of intrinsic motivation)
		\end{itemize}
		
		\vspace{0.3cm}
		
		
	\end{itemize}
	
\end{frame}

% --- Research gap \& Aim of the paper ---
\begin{frame}{Research Questions} % 
	\begin{itemize}
		\item This study focuses on intrinsic motivation and evaluates interventions
		that can improve the effectiveness of extension activities.
		
		\vspace{0.3cm}
		
		\begin{itemize}
			
			\item \textbf{Is low extension effort driven by a lack of locus of control?}
			\begin{itemize}
				\item Extension workers do not believe that their efforts can improve farmers' lives.
			\end{itemize}
			
			\vspace{0.3cm}
			
			\item \textbf{Is low effort driven by weak social recognition?}
			\begin{itemize}
				\item No matter how hard AEAs try, their efforts are not particularly recognized or appreciated by their superiors or the farmers. 
			\end{itemize}
			
			\vspace{0.3cm}
			
			
		\end{itemize}
		
	\end{itemize}
\end{frame}

\begin{frame}{What We Do}
	
	\begin{itemize}
		\item This study focuses on intrinsic motivation and evaluates interventions
		that can improve the effectiveness of extension activities.
		
		\vspace{0.3cm}
		\begin{itemize}
			\item \textbf{Is low extension effort driven by a lack of locus of control?}
			\begin{itemize}
				\item Extension workers do not believe that their efforts can improve farmers' lives.
				\item $\Rightarrow$ Provide training emphasizing the social significance of extension activities
			\end{itemize}
			
			\vspace{0.3cm}
			
			\item \textbf{Is low effort driven by weak social recognition?}
			\begin{itemize}
				\item  No matter how hard AEAs try, their efforts are not particularly recognized or appreciated by their superiors or the farmers. 
				\item $\Rightarrow$ Introduce a feedback system communicating farmers' satisfaction
				\item $\Rightarrow$ Ask AEAs to stamp their location information when visit farmer fields (somewhat related to extrinsic motivation?)
			\end{itemize}
			
			\vspace{0.3cm}
			
			
		\end{itemize}
	\end{itemize}
	
\end{frame}

\begin{frame}{Study Areas}
	\begin{itemize}
		\item Rainfed areas covered by GRIP:
		\begin{itemize}
			\item Oti (Biakoye and Krachi West)
			\item Bono East (Pru West and Sene West)
			\item Western North (Bibiani and Sefwi Wiaoso)
			\item Eastern Regions (Achiase and Lower Manya) 
			\vspace{0.3cm}
		\end{itemize}
		\item Irrigation schemes:
		\begin{itemize}
			\item Weta (Volta)
			\item Kpong (Eastern)
		\end{itemize}
		\vspace{0.3cm}
		\item Total: 
		\begin{itemize}
			\item 10 study sites across major rice-producing environments in Ghana
			\item 44 agricultural extension agents (AEAs)
			\item 390 rice-growing farmers
			
		\end{itemize}
		
	\end{itemize}
\end{frame}

\begin{frame}{Random Assignment}
	\textbf{Sample}
	\begin{itemize}
		\item 44 AEAs who participated in the baseline survey
		\item Random assignment conducted at the district level
	\end{itemize}
	
	
	\textbf{Treatment Arms}
	\begin{itemize}
		\item \textbf{Group 1: Feedback Treatment}
		\begin{itemize}
			\item Extension agents and their supervisors receive feedback from farmers regarding their extension activities: biweekly
			\item \textcolor{blue}{Sene West, Achiase, Bibiani}
		\end{itemize}
		
		\item \textbf{Group 2: Feedback + Training}
		\begin{itemize}
			\item Farmer feedback provided
			\item Additional training emphasizing the social importance and impact of extension services
			\item \textcolor{blue}{Lower Manya, Krachi West, Sefwi Wiaoso}
		\end{itemize}
		
		\item \textbf{Group 3: Control Group}
		\begin{itemize}
			\item No intervention (Collect the satisfaction level, but do not inform the result to AEAs) 
			\item \textcolor{blue}{Pru West, Biakoye, and 2 irrigation schemes}
		\end{itemize}
		
		\item \textbf{Half of AEAs are subject to stamp treatment}
		
	\end{itemize}
	
\end{frame}

\begin{frame}{Main outcomes of interest}
	\begin{itemize}
		\item Rice Production (Farmer data at baseline and endline)
		\begin{itemize}
			\item Yield
			\item Income 
			\item Profits
			\item Rice management practices
			\vspace{0.3cm}
		\end{itemize}
		\item Extension Activities (Biweekly data) 
		\begin{itemize}
			\item Number of visits
			\item Number of successful communication
			\item Satisfaction level 
		\end{itemize}
		
		
		
	\end{itemize}
\end{frame}

\begin{frame}{Project Timeline}
	
	\begin{itemize}
		\item \textbf{January 2025:} Baseline Survey
		\vspace{0.3cm}
		
		\item \textbf{April -- September 2025:} Intervention
		\begin{itemize}
			\item Satisfaction Survey
			\item Training
		\end{itemize}
		\vspace{0.3cm}
		
		\item \textbf{January 2026:} Endline Survey
	\end{itemize}
	
\end{frame}

\begin{frame}{Number of Obs by District}
\begin{table}[hbtp]
	\centering
	\begin{threeparttable}
		
		% ここで列の定義を自分で行う
		\begin{tabular}{llccc}
			\hline
			\input{tmp/obs_data.tex} % Stataが出力した中身だけを読み込む
			
		\end{tabular}
		
	\end{threeparttable}
\end{table}
\end{frame}	




\begin{frame}{Demographic Characteristics (Farmer)}
	\centering
	\scalebox{0.55}{\input{tmp/f_balance1.tex}}
	
\end{frame}

\begin{frame}{Basic Farming Characteristics (Farmer)}
	\centering
	\scalebox{0.5}{\input{tmp/f_balance2.tex}}
	
\end{frame}


\begin{frame}{Basic Farming Characteristics (Farmer) w/o Irrg}
\centering
\scalebox{0.5}{\input{tmp/f_balance2b.tex}}

\end{frame}

\begin{frame}{Technology Adoption (Farmer)}
	\centering
	\scalebox{0.6}{\input{tmp/f_balance3.tex}}
	
\end{frame}

\begin{frame}{Basic Characteristics (AEA)}
	\centering
	\scalebox{0.55}{\input{tmp/a_balance1.tex}}
	
\end{frame}

\begin{frame}{Prosocialness (AEA)}
	\centering
	\scalebox{0.6}{\input{tmp/a_balance2.tex}}
	
\end{frame}

\begin{frame}{Intrinsic and Extrinsic Motivation}
	\begin{itemize}
		\item Intrinsic motivation
		\begin{itemize}
			\item Because this activity is fun 
			\item Because I feel good when doing this activity
			\vspace{0.3cm}
		\end{itemize}
		\item Extrinsic Motivation (External regulation and identified regulation)
		\begin{itemize}
			\item Because I am supposed to do it
			\item Because I don't have any choice
			\item Because I am doing it for my own good 
			\item Because I think that this activity is good for me. 
		\end{itemize}
		
		
		
	\end{itemize}
\end{frame}


\begin{frame}{Distribution of Prosocialness Answer}
	\begin{figure}
		\includegraphics[width=12cm]{tmp/prosocial.eps}
	\end{figure}
\end{frame}

\begin{frame}{Locus of Control (AEA)}
	\centering
	\scalebox{0.5}{\input{tmp/a_balance3.tex}}
	
\end{frame}

\begin{frame}{Distribution of Locus of Control Answer (A-E)}
	\begin{figure}
		\includegraphics[width=12cm]{tmp/locus1.eps}
	\end{figure}
\end{frame}

\begin{frame}{Distribution of Locus of Control Answer (F-J)}
	\begin{figure}
		\includegraphics[width=12cm]{tmp/locus2.eps}
	\end{figure}
\end{frame}

\begin{frame}{Descriptive Results of Biweekly Survey}
	\centering
\scalebox{0.6}{\input{tmp/b_balance1.tex}}
\end{frame}




\begin{frame}{Descriptive Results of Biweekly Survey}
	\begin{figure}
		\includegraphics[width=12cm]{tmp/b_basic.eps}
	\end{figure}
\end{frame}

\begin{frame}{Satisfaction Level by District and Time}
	\begin{figure}
		\includegraphics[width=10cm]{tmp/b_satisfy.eps}
	\end{figure}
\end{frame}


\begin{frame}{Preliminary results (AEA)}
	\centering
	\scalebox{0.55}{\input{tmp/result_table1.tex}}
		\begin{tablenotes}
			\scriptsize
		\item Note: Provincial and Round fixed effects are controlled. ***, **, and * indicate significance at the 1, 5, and 10 percent critical level.  \\.  
	\end{tablenotes}
	
\end{frame}

\begin{frame}{Preliminary results (AEA)}
	\centering
	\scalebox{0.55}{\input{tmp/result_table2.tex}}
	\begin{tablenotes}
		\scriptsize
		\item Note: Provincial and Round fixed effects are controlled. ***, **, and * indicate significance at the 1, 5, and 10 percent critical level.  \\.  
	\end{tablenotes}
	
\end{frame}

\begin{frame}{Preliminary results of Correlation (Farming)}
	\centering
	\scalebox{0.34}{\input{tmp/f_result_table2.tex}}
	
\end{frame}

\begin{frame}{Preliminary results of Correlation (AEA)}
	\centering
	\scalebox{0.5}{\input{tmp/a_result_table3.tex}}
	
\end{frame}

\begin{frame}{Ways forward}
\begin{itemize}
	\item More data cleaning
	\item Endline survey (January-March/April)
	\item Analysis 
	\begin{itemize}
		\item Using sampling weight
		\item Wild bootstrap
		\item Randomization inference 
	\end{itemize}
	\item Follow-up survey? 
\end{itemize}
	
\end{frame}
	
\end{document}
